% Generated by tools/edn_latex.py from reports/edn.md.
% Do not edit: change reports/edn.md and run `make edn`.
\ednentry{
  id = {EDN-58},
  title = {An unmeasured success criterion blocks the gate, and the artefact says how close the model is anyway},
  date = {2026-09-30},
  milestone = {M3: Quality Assurance},
  topic = {Testing Strategy, Model Evaluation (the NFR-01 gate)},
  participants = {@lukas2510},
  decision = {A criterion the \texttt{evaluate} stage could not measure is recorded as \texttt{null}, never as a pass and never as a failure, and \texttt{null} blocks NFR-01's gate. Two criteria are in that state today: SC-05, because no model exposes \texttt{predict\_\allowbreak{}interval\_\allowbreak{}eur} and the UC2 conformal intervals are a later ticket, and SC-04 whenever no segment level reaches the 500-row minimum. Alongside the blocking verdict, \texttt{metrics.\allowbreak{}json} carries two leaves that say what is outstanding and how close the candidates are without it: \texttt{criteria\_\allowbreak{}not\_\allowbreak{}measured} names the criteria, and \texttt{n\_\allowbreak{}variants\_\allowbreak{}passing\_\allowbreak{}measurable} counts the variants where nothing that \emph{could} be measured failed. Each criterion also carries a structured status rather than only prose: \texttt{sc05\_\allowbreak{}status: \textquotedbl{}not\_\allowbreak{}measured\textquotedbl{}}, \texttt{sc05\_\allowbreak{}capability\_\allowbreak{}checked: \textquotedbl{}predict\_\allowbreak{}interval\_\allowbreak{}eur\textquotedbl{}}, \texttt{sc05\_\allowbreak{}reason: \textquotedbl{}.\allowbreak{}.\allowbreak{}.\allowbreak{}\textquotedbl{}}, and \texttt{sc04\_\allowbreak{}note} naming the largest level and its row count.},
  rationale = {\emph{Alternatives considered:}
\begin{itemize}[leftmargin=*, topsep=2pt]
\item \textbf{Option A (chosen): block, with a structured \texttt{null} and a second, clearly named verdict.}
\newline Pros: NFR-01's word is ``every'', and this is the only reading of it that survives an unbuilt criterion. A reader of the artefact can tell ``we did not measure it'' from ``the model missed it'' without parsing English, which is what EDN-12 needs, since promotion is a human decision and the human has to see what is outstanding. \texttt{n\_\allowbreak{}variants\_\allowbreak{}passing\_\allowbreak{}measurable} keeps the first delivery's metrics file informative: \texttt{n\_\allowbreak{}variants\_\allowbreak{}passing: 0} says nothing about the model, and this says how close the ladder is.
\newline Cons: two leaves and a status string per criterion that a one-criterion-at-a-time reader does not need, and a second verdict that a careless reader could quote as if it were the gate.
\item \textbf{Option B: block, with \texttt{null} and nothing else.}
\newline Pros: the narrowest possible \texttt{metrics.\allowbreak{}json}, which is what the existing comment in \texttt{dvc.\allowbreak{}yaml} was written to defend.
\newline Cons: the first delivery's artefact then reads ``0 of 4 variants pass'' with no explanation anywhere in it, and the reason is a ticket that has not been written rather than anything about the models. That is the number the report would have to cite, and it would have to be explained in prose beside the artefact instead of in it.
\item \textbf{Option C: exclude an unmeasurable criterion from the gate and pass on the rest.}
\newline Pros: the gate would produce a deployable candidate now, which is what a demo wants.
\newline Cons: unacceptable. It would let a model pass NFR-01 without any interval ever being calibrated, which is precisely the promise the requirement makes. It also gets easier to do the more criteria are outstanding, so the gate would weaken exactly when it matters most.
\item \textbf{Option D (for SC-04 specifically): treat ``every level satisfies P'' as met when no level qualifies.}
\newline Pros: vacuously true, and it is what \texttt{all([])} returns, so it is what the obvious implementation does.
\newline Cons: it reports a check nobody ran as a check that passed. On the 2,000-row test fixture no level reaches 500 rows, so this option would have reported SC-04 as met for every variant in every CI run.
\end{itemize}
\par
\emph{Why:} Option C is the only one that is actually wrong, and the choice between A and B is about whether the artefact has to be readable on its own. It does: EDN-12 makes promotion a pull-request review, and the thing the reviewer reads is \texttt{dvc metrics diff}, which is long-format and lists only the leaves that changed. A row reading \texttt{criteria\_\allowbreak{}not\_\allowbreak{}measured: sc05 -\allowbreak{}> \textquotedbl{}\textquotedbl{}} is the single most useful line that diff can carry, and it does not exist under option B. Option D was rejected on a measurement rather than on taste: at the committed fixture size zero segment levels clear the 500-row bar, so the vacuous reading would have been the one CI exercised, and an implementation of SC-04 that simply returned ``passed'' would have kept the suite green. The test suite answers that by running \texttt{evaluate} a second time with \texttt{sc04\_\allowbreak{}min\_\allowbreak{}segment\_\allowbreak{}rows} lowered to 50, which produces a real verdict on the fixture; growing \texttt{FIXTURE\_\allowbreak{}ROWS} from 2,000 to about 20,000 so the project's own bar is exercised was considered and rejected, because it multiplies the cost of the session fixture across the whole suite for a property the parameterised run already covers more precisely, and the fixture size is itself a recorded decision (EDN-33).},
  ai = {Alternative generation, Alternative assessment, Recommendation, Solution generation},
  response = {Accepted with modifications},
  assessment = {AI laid out the four readings of an unmeasurable criterion and recommended A, with the argument that option C weakens the gate more the more criteria are outstanding. The part that was worth more than the recommendation was the measurement behind option D: AI ran the split and the segment cut at three data scales before writing any of the stage, found that no level clears 500 rows on the pytest fixture, and concluded that the vacuous reading would be the one CI exercised and that a test asserting \texttt{sc04\_\allowbreak{}passed is None} there would be satisfied by an implementation that returned \texttt{None} unconditionally. That is the trap the second \texttt{evaluate} run in \texttt{tests/\allowbreak{}test\_\allowbreak{}pipeline.\allowbreak{}py} exists for, and it came out of the measurement rather than out of the design. The modification was to the leaf count: an earlier proposal also put \texttt{sc04\_\allowbreak{}worst\_\allowbreak{}segment} and the per-field sweep into \texttt{metrics.\allowbreak{}json}, which was moved to \texttt{reports/\allowbreak{}metrics/\allowbreak{}<variant>.\allowbreak{}json} to keep \texttt{dvc metrics show} at 24 columns.},
  aievidence = {Claude Code sessions on 2026-09-30 designing and then implementing issue \#39. The design session measured the split and segment sizes at 416, 3,695 and 22,551 test rows and reported ``SC-04's 500-row rule is not measurable at all on the 2,000-row pytest fixture; zero levels qualify, the largest is \texttt{seller\_\allowbreak{}type=\allowbreak{}Dealer} at 345 rows''; it also validated the proposed \texttt{metrics.\allowbreak{}json} shape against \texttt{dvc metrics show} and \texttt{dvc metrics diff} in a throwaway DVC repository before the shape was fixed.},
  evidence = {\href{https://github.com/mlops-2627q1-mds-upc/MLOPS_Recommenditos/blob/main/recommenditos/modeling/evaluate.py}{\texttt{recommenditos/\allowbreak{}modeling/\allowbreak{}evaluate.\allowbreak{}py}}, \texttt{evaluate\_\allowbreak{}sc04}, \texttt{evaluate\_\allowbreak{}sc05}, \texttt{gate\_\allowbreak{}passed}, \texttt{gate\_\allowbreak{}blocked\_\allowbreak{}by} and \texttt{\_\allowbreak{}check\_\allowbreak{}every\_\allowbreak{}criterion\_\allowbreak{}is\_\allowbreak{}decided}; \texttt{tests/\allowbreak{}test\_\allowbreak{}model.\allowbreak{}py::test\_\allowbreak{}sc04\_\allowbreak{}is\_\allowbreak{}not\_\allowbreak{}measured\_\allowbreak{}rather\_\allowbreak{}than\_\allowbreak{}vacuously\_\allowbreak{}met\_\allowbreak{}when\_\allowbreak{}no\_\allowbreak{}level\_\allowbreak{}qualifies}, \texttt{::test\_\allowbreak{}an\_\allowbreak{}unmeasured\_\allowbreak{}criterion\_\allowbreak{}blocks\_\allowbreak{}the\_\allowbreak{}gate\_\allowbreak{}rather\_\allowbreak{}than\_\allowbreak{}passing\_\allowbreak{}it}, \texttt{::test\_\allowbreak{}a\_\allowbreak{}criterion\_\allowbreak{}added\_\allowbreak{}to\_\allowbreak{}the\_\allowbreak{}list\_\allowbreak{}without\_\allowbreak{}a\_\allowbreak{}measurement\_\allowbreak{}fails\_\allowbreak{}the\_\allowbreak{}stage}; \texttt{tests/\allowbreak{}test\_\allowbreak{}pipeline.\allowbreak{}py::test\_\allowbreak{}sc04\_\allowbreak{}is\_\allowbreak{}measured\_\allowbreak{}as\_\allowbreak{}soon\_\allowbreak{}as\_\allowbreak{}a\_\allowbreak{}segment\_\allowbreak{}level\_\allowbreak{}is\_\allowbreak{}large\_\allowbreak{}enough}; \href{https://github.com/mlops-2627q1-mds-upc/MLOPS_Recommenditos/blob/main/docs/docs/problem-spec.md}{problem specification} section 8; \hyperref[EDN-06]{EDN-06}, \hyperref[EDN-12]{EDN-12}, \hyperref[EDN-33]{EDN-33}; \href{https://github.com/mlops-2627q1-mds-upc/MLOPS_Recommenditos/issues/39}{issue \#39}.},
}
